\documentclass[11pt]{article}
\usepackage[margin=1in]{geometry}
\usepackage[T1]{fontenc}
\usepackage{lmodern}
\usepackage{amsmath,amssymb,amsthm}
\usepackage[hidelinks]{hyperref}
\usepackage{microtype}

\newtheorem{theorem}{Theorem}
\newtheorem{lemma}[theorem]{Lemma}
\newcommand{\R}{\mathbb R}
\newcommand{\G}{\mathcal G}

\begin{document}

\begin{center}
{\Large Full-result packet for arXiv:2603.11040}\\[4pt]
{\large Exact one-point faithfulness and classification of all optimizers}
\end{center}

\paragraph{Status.}
\textbf{full\_solution\_likely\_valid}.  This packet sharpens the one-point
thresholding theorem in Sujit Sakharam Damase and James Eldred Pascoe,
\emph{On positive definite thresholding of correlation matrices},
arXiv:2603.11040v1.

\paragraph{Source problem.}
Fix \(n\geq3\), put \(\alpha=(n-2)/2\), and write
\[
  \G_k(t)=\frac{C_k^{(\alpha)}(t)}{C_k^{(\alpha)}(1)}
\]
for the normalized Gegenbauer polynomial, so \(\G_k(1)=1\) and
\(\G_1(t)=t\).  Schoenberg's theorem says that the continuous unital
positive definite functions on \(S^{n-1}\) are exactly
\[
  f(t)=\sum_{k=0}^{\infty}a_k\G_k(t),
  \qquad a_k\geq0,\qquad \sum_{k=0}^{\infty}a_k=1.
\]
For \(K\subset[-1,1)\), the source defines the faithfulness constant
\(\tau_{K,n}\) as the largest possible linear coefficient \(a_1\) among
such functions vanishing on \(K\).  For \(K=\{\varepsilon\}\), with
\(\varepsilon>0\) small, the source constructs a degree-two function and
proves only
\[
 \tau_{\{\varepsilon\},n}\geq
 \frac{|\G_2(\varepsilon)|}
      {\varepsilon+|\G_2(\varepsilon)|}.
\]
The question addressed here is the exact value and the equality case.

\paragraph{Result.}
The one-point problem is an infinite-dimensional linear program with only
one interpolation constraint.  Its optimum is determined by the most
negative Gegenbauer value at the threshold point.  Moreover, every optimizer
is supported on degree one and the degrees attaining that minimum.  A
uniform Darboux estimate then shows that, for every fixed \(n\geq3\), degree
two is the unique minimizing degree for all sufficiently small positive
thresholds.  Thus the source's lower bound is locally an equality, and its
quadratic witness is the unique optimizer.

\begin{lemma}\label{lem:min}
For every \(n\geq3\) and \(\varepsilon\in(0,1)\), the number
\[
  m_n(\varepsilon)
  :=\min_{k\in\mathbb Z_{\geq0}\setminus\{1\}}\G_k(\varepsilon)
\]
is well defined and is strictly negative.
\end{lemma}

\begin{proof}
The existence theorem in arXiv:2603.11040 supplies a nonzero continuous
positive definite function vanishing at \(\varepsilon\).  After normalizing
at \(1\), its Schoenberg coefficients form a probability vector.  Since
\(\G_0(\varepsilon)=1\) and \(\G_1(\varepsilon)=\varepsilon>0\), the identity
\(\sum_k a_k\G_k(\varepsilon)=0\) forces
\(\G_j(\varepsilon)<0\) for some \(j\ne1\).

For fixed \(\varepsilon\in(-1,1)\), the classical Darboux estimate for
normalized Gegenbauer polynomials gives
\(\G_k(\varepsilon)\to0\) as \(k\to\infty\) when \(n\geq3\).  Because at
least one term is negative, the infimum is negative and is attained among
finitely many indices.  It is therefore the displayed minimum.
\end{proof}

\begin{theorem}[Exact one-point faithfulness]\label{thm:exact}
Let \(n\geq3\) and \(0<\varepsilon<1\).  Then
\[
 \boxed{\displaystyle
 \tau_{\{\varepsilon\},n}
 =\frac{-m_n(\varepsilon)}
        {\varepsilon-m_n(\varepsilon)}}.
\]
Let
\(M_n(\varepsilon)=\{k\ne1:\G_k(\varepsilon)=m_n(\varepsilon)\}\).
A Schoenberg coefficient vector \(a=(a_k)\) is optimal if and only if
\[
 a_1=\frac{-m_n(\varepsilon)}
           {\varepsilon-m_n(\varepsilon)},
 \qquad
 \sum_{k\in M_n(\varepsilon)}a_k
   =\frac{\varepsilon}{\varepsilon-m_n(\varepsilon)},
\]
and every other coefficient vanishes.  In particular, an optimizer can
always be chosen to use exactly two Gegenbauer modes.
\end{theorem}

\begin{proof}
Let \(f=\sum a_k\G_k\) be feasible.  Using \(f(\varepsilon)=0\), the
definition of \(m=m_n(\varepsilon)\), and \(\sum_k a_k=1\), we obtain
\[
 0=a_1\varepsilon+\sum_{k\ne1}a_k\G_k(\varepsilon)
 \geq a_1\varepsilon+m\sum_{k\ne1}a_k
 =a_1\varepsilon+m(1-a_1).
\]
As \(m<0\), this rearranges to
\[
 a_1\leq\frac{-m}{\varepsilon-m}.
\]
Choose \(j\in M_n(\varepsilon)\) and set
\[
 a_1=\frac{-m}{\varepsilon-m},
 \qquad
 a_j=\frac{\varepsilon}{\varepsilon-m},
 \qquad a_k=0\quad(k\notin\{1,j\}).
\]
These coefficients are nonnegative, sum to one, and satisfy
\(a_1\varepsilon+a_jm=0\).  They therefore define a feasible continuous
positive definite function attaining the upper bound.

Finally, equality in the displayed estimate is possible precisely when
every positive coefficient outside degree one is carried by an index for
which \(\G_k(\varepsilon)=m\).  This is exactly the asserted classification.
\end{proof}

\begin{theorem}[Local sharpness and uniqueness of the quadratic witness]
For every fixed \(n\geq3\), there is \(\delta_n>0\) such that, whenever
\(0<\varepsilon<\delta_n\),
\[
 m_n(\varepsilon)=\G_2(\varepsilon)
 =\frac{n\varepsilon^2-1}{n-1},
\]
and degree two is the unique minimizing degree.  Consequently
\[
 \boxed{\displaystyle
 \tau_{\{\varepsilon\},n}
 =\frac{1-n\varepsilon^2}
 {1-n\varepsilon^2+(n-1)\varepsilon}}
\]
for all such \(\varepsilon\), and the unique optimal function is
\[
 f_\varepsilon(t)
 =\tau_{\{\varepsilon\},n}\,t
 +(1-\tau_{\{\varepsilon\},n})
   \frac{nt^2-1}{n-1}.
\]
\end{theorem}

\begin{proof}
At zero, parity gives \(\G_{2r+1}(0)=0\), while the duplication formula for
the Gamma function gives
\[
 \G_{2r}(0)
 =(-1)^r\frac{(1/2)_r}{((n-1)/2)_r}
 \qquad(r\geq0).
\]
For \(n\geq3\), the absolute values on the right strictly decrease with
\(r\).  Hence
\[
 \G_2(0)=-\frac1{n-1}
 <\G_k(0)\qquad(k\ne1,2),
\]
with a positive gap: the even subsequence decreases in absolute value and
the odd values are zero.

The Darboux estimate is uniform for \(t\) in a compact subinterval of
\((-1,1)\), so \(\G_k(t)\to0\) uniformly for \(t\) in a fixed neighborhood
of zero.  Choose the tail large enough that no tail polynomial can close the
gap to \(\G_2(0)\).  The remaining family is finite, and continuity preserves
the strict inequalities after shrinking the neighborhood.  Thus there is
\(\delta_n>0\) for which degree two remains the unique minimizer.

The formula
\(\G_2(t)=(nt^2-1)/(n-1)\), followed by Theorem~\ref{thm:exact}, gives the
displayed value and the unique equality case.
\end{proof}

\paragraph{Checks and scope.}
The proof uses exactly the source's continuous Schoenberg class and the same
normalization of the faithfulness coefficient.  It is restricted to
\(n\geq3\), where normalized Gegenbauer values at an interior point tend to
zero; the circle \(n=2\) has a distinct Chebyshev/Diophantine endpoint and is
not claimed here.  The result sharpens the source's one-point theorem but
does not resolve its final conjecture about discontinuous positive definite
functions.

\paragraph{Novelty check.}
Before promotion, the run indexes were searched for arXiv:2603.11040, the
phrases ``faithfulness constant'' and ``one-point thresholding,'' and the
Gegenbauer/positive-definite combinations; no earlier packet was found.  A
bounded primary-source search on 2026-07-19 found only arXiv:2603.11040v1 and
its companion arXiv:2602.09010.  The current source record still describes
``bounds'' at single points, and v1 contains only the quadratic lower bound,
not the exact formula or equality classification above.

\paragraph{Human-review recommendation.}
Review as a short exact extension of the source's one-point theorem.  The
linear-program argument is elementary; the only analytic point to audit is
the standard uniform-on-compacts Darboux estimate used to make degree two
uniformly optimal near zero.

\begin{thebibliography}{9}
\bibitem{source}
S. S. Damase and J. E. Pascoe,
\emph{On positive definite thresholding of correlation matrices},
arXiv:2603.11040v1, 2026.

\bibitem{companion}
S. S. Damase and J. E. Pascoe,
\emph{Complete discrete Schoenberg--Delsarte theory for homogeneous spaces},
arXiv:2602.09010v1, 2026.

\bibitem{schoenberg}
I. J. Schoenberg,
\emph{Positive definite functions on spheres},
Duke Math. J. \textbf{9} (1942), 96--108.
\end{thebibliography}

\end{document}
